\begin{textsample}{POS Dim 4 – human – Score 6.00 – mg/mg\_ef\_linguagens\_ingles\_8\_p8.txt}  \label{ex:f4_pos_001}
A competência \textbf{textual} inclui, por exemplo, a capacidade do estudante de saber distinguir um bilhete de um anúncio publicitário, um horóscopo de uma receita culinária e de saber verificar se, em um determinado \textbf{texto}, predominam sequências linguísticas de caráter narrativo, descritivo ou argumentativo.
Claro está que o estudante será mais ( ou menos ) capaz de estabelecer essas distinções e de reconhecer os diferentes padrões de organização \textbf{textual}, dependendo do seu nível de escolaridade e de suas \textbf{oportunidades} de contato com \textbf{textos} \textbf{orais} ou escritos no meio em que vive. É certo que o estudante do Ensino Fundamental e Médio possui conhecimento genérico sobre as situações e eventos habituais do dia a dia armazenados na memória de longo prazo ( conhecimento prévio ou enciclopédico ).
Diante disso, parece possível afirmar que ele também possui, pelo fato de viver numa sociedade altamente letrada, informações amplas sobre os vários gêneros \textbf{textuais} como cartas, bilhetes, guias turísticos, anúncios publicitários, etc.

% matched lemmas: oportunidade, oral, texto, textual
\end{textsample}
